\chapter{HTTP: messaggi, cookie e caching}

Nel capitolo precedente abbiamo visto \emph{come} HTTP usa le connessioni. Ricapitoliamo e passiamo
a \emph{che cosa} si scambiano client e server.

\gbox{Ripasso: HTTP in cinque punti}{primary!70!black}{
\begin{itemize}
    \item È uno dei protocolli più usati di Internet (applicazioni Web).
    \item È \textbf{client-server} e definisce \textbf{il formato dei messaggi} e \textbf{come vengono
          scambiati} (non come si implementano client e server).
    \item Usa messaggi di \textbf{richiesta} e di \textbf{risposta}.
    \item È \textbf{stateless}: client e server non ricordano le interazioni passate.
    \item Si appoggia a un trasporto affidabile e orientato alla connessione (soprattutto TCP, di recente
          anche UDP con QUIC), con connessioni persistenti o non persistenti.
\end{itemize}
}

% ══════════════════════════════════════════════════════════════
\section{Il formato dei messaggi HTTP}
% ══════════════════════════════════════════════════════════════

HTTP ha due formati diversi per i due tipi di messaggio:

\begin{itemize}
    \item la \textbf{richiesta} specifica il \textbf{comando} (metodo) che il server deve eseguire: dammi
          una pagina, compila questo modulo, \dots;
    \item la \textbf{risposta} riporta l'\textbf{esito} del comando (successo, errore, \dots) ed eventuali
          \textbf{dati} (per esempio il file richiesto).
\end{itemize}

Entrambi sono scritti in \textbf{testo ASCII} ordinario: si leggono facilmente anche a occhio.

\subsection{Il formato della richiesta}

\begin{figure}[H]
\centering
\begin{tikzpicture}[every node/.style={font=\scriptsize\sffamily}, f/.style={field, minimum height=0.6cm}]
  % request line
  \node[f, fill=lapp, minimum width=1.6cm] (m) at (0,0) {metodo};
  \node[f, fill=white, minimum width=0.5cm, right=0pt of m] (s1) {sp};
  \node[f, fill=lnet, minimum width=3cm, right=0pt of s1] (u) {URL};
  \node[f, fill=white, minimum width=0.5cm, right=0pt of u] (s2) {sp};
  \node[f, fill=ltra, minimum width=1.6cm, right=0pt of s2] (v) {versione};
  \node[f, fill=white, minimum width=0.5cm, right=0pt of v] {cr}; \node[f, fill=white, minimum width=0.5cm] at ([xshift=0.75cm]v.east) {lf};
  % header lines
  \foreach \y in {-0.8,-1.9} {
    \node[f, fill=llink, minimum width=3.1cm, anchor=west] (h) at (-0.8,\y) {nome campo header:};
    \node[f, fill=white, minimum width=0.5cm, right=0pt of h] (sp) {sp};
    \node[f, fill=llink!50, minimum width=2.2cm, right=0pt of sp] (val) {valore};
    \node[f, fill=white, minimum width=0.5cm, right=0pt of val] (cr) {cr};
    \node[f, fill=white, minimum width=0.5cm, right=0pt of cr] {lf};
  }
  \node at (1.5,-1.35) {$\vdots$};
  \node[f, fill=white, minimum width=0.5cm, anchor=west] (bcr) at (-0.8,-2.7) {cr};
  \node[f, fill=white, minimum width=0.5cm, right=0pt of bcr] {lf};
  \node[f, fill=lphy!60, minimum width=8.8cm, minimum height=1cm, anchor=north west] at (-0.8,-3.1) {corpo (\emph{entity body}): dati, se il metodo li prevede};
  % labels
  \draw[decorate, decoration={brace, amplitude=5pt}] (-1.1,-0.3) -- (-1.1,0.3) node[midway, left=6pt] {riga di richiesta};
  \draw[decorate, decoration={brace, amplitude=5pt}] (-1.1,-2.2) -- (-1.1,-0.5) node[midway, left=6pt] {righe di header};
  \node[left] at (-1.3,-2.7) {riga vuota};
  \node[left] at (-1.3,-3.6) {corpo};
\end{tikzpicture}
\caption{Formato generale di un messaggio di richiesta HTTP.}
\end{figure}

Il messaggio di richiesta contiene:

\begin{itemize}
    \item il \textbf{metodo}: il comando da far eseguire al server;
    \item l'\textbf{URL}: l'oggetto su cui si vuole operare;
    \item la \textbf{versione} di HTTP (per esempio \texttt{HTTP/1.1});
    \item le \textbf{righe di header}: i parametri della richiesta, in numero e tipo non fissati; ogni riga ha
          il nome e il valore del parametro. Qui si può per esempio chiedere una connessione persistente o
          no; si possono usare anche header personalizzati;
    \item il \textbf{corpo} (\emph{body}): dipende dal metodo e contiene gli eventuali dati associati al
          comando (per esempio il testo con cui compilare un modulo).
\end{itemize}

\warning{I separatori}{
    I campi sono separati da caratteri speciali: \texttt{sp} è lo \textbf{spazio}; \texttt{cr} è il
    \emph{carriage return} (\texttt{\textbackslash r}); \texttt{lf} è il \emph{line feed}
    (\texttt{\textbackslash n}). Ogni riga termina con \texttt{\textbackslash r\textbackslash n}, e una
    \textbf{riga vuota} (\texttt{\textbackslash r\textbackslash n} da solo) separa gli header dal corpo.
    Dimenticarla, quando si scrive una richiesta a mano, è l'errore più comune.
}

\subsection{I metodi HTTP}

\begin{table}[H]
\centering
\begin{tabular}{lp{12cm}}
\toprule
\textbf{Metodo} & \textbf{Uso} \\
\midrule
\texttt{GET} & \textbf{recupera} un oggetto dal server. È il metodo più usato: ogni volta che si apre una pagina si
fa una GET. Il corpo è \textbf{vuoto}. \\
\texttt{POST} & \textbf{invia informazioni} al server, per esempio il contenuto di un modulo. Il corpo contiene
i dati. \\
\texttt{HEAD} & come GET, ma il server risponde \textbf{senza l'oggetto}, con i soli header. Utile per il
debugging e per controllare se un oggetto esiste o è cambiato. \\
\texttt{PUT} & \textbf{carica} un oggetto in un certo percorso del server (strumenti di pubblicazione Web,
applicazioni che devono caricare file). \\
\texttt{DELETE} & \textbf{cancella} un oggetto dal server. \\
\bottomrule
\end{tabular}
\caption{I metodi HTTP principali.}
\end{table}

\warning{I moduli non usano per forza POST}{
    Una richiesta generata da un modulo HTML non usa necessariamente POST: spesso i moduli usano GET e
    inseriscono i dati dei campi \textbf{direttamente nell'URL}, come query (per esempio
    \texttt{...search?q=reti\&lang=it}).
}

\subsection{Esempio di richiesta}

\begin{lstlisting}[language=http, title={Richiesta HTTP}]
GET /somedir/page.html HTTP/1.1
Host: www.someschool.edu
Connection: close
User-agent: Mozilla/5.0
Accept-language: fr
\end{lstlisting}

Un browser (Mozilla/5.0) che implementa HTTP/1.1 chiede l'oggetto \texttt{/somedir/page.html} al server
\texttt{www.someschool.edu}. Ci sono cinque righe: la riga di richiesta (GET, risorsa, versione) e quattro
righe di header; il corpo è vuoto.

\begin{itemize}
    \item \texttt{Host: www.someschool.edu} indica l'host su cui si trova l'oggetto. Sembra ridondante (ci siamo
          già connessi a quell'host!), ma il destinatario della connessione non è sempre il server finale: il
          messaggio può essere inoltrato da un \textbf{proxy}, e allora l'indirizzo della connessione TCP è
          diverso da quello del server vero.
    \item \texttt{Connection: close} chiede una connessione \textbf{non persistente}: chiudi quando hai finito.
    \item \texttt{User-agent: Mozilla/5.0} indica il tipo di browser: il server può inviare versioni diverse
          dello stesso oggetto (stesso URL) a seconda del client.
    \item \texttt{Accept-language: fr} indica che l'utente preferisce, se esiste, la versione in francese.
\end{itemize}

\subsection{Il formato della risposta}

La risposta ha un formato analogo, ma al posto della riga di richiesta c'è una \textbf{riga di stato}
che riporta l'esito del comando:

\begin{figure}[H]
\centering
\begin{tikzpicture}[every node/.style={font=\scriptsize\sffamily}, f/.style={field, minimum height=0.6cm}]
  \node[f, fill=ltra, minimum width=1.6cm] (v) at (0,0) {versione};
  \node[f, fill=white, minimum width=0.5cm, right=0pt of v] (s1) {sp};
  \node[f, fill=lapp, minimum width=2cm, right=0pt of s1] (c) {codice di stato};
  \node[f, fill=white, minimum width=0.5cm, right=0pt of c] (s2) {sp};
  \node[f, fill=lnet, minimum width=1.8cm, right=0pt of s2] (p) {frase};
  \node[f, fill=white, minimum width=0.5cm, right=0pt of p] (cr) {cr};
  \node[f, fill=white, minimum width=0.5cm, right=0pt of cr] {lf};
  \node[f, fill=llink, minimum width=7.4cm, anchor=west] at (-0.8,-0.8) {righe di header (come nella richiesta)};
  \node[f, fill=white, minimum width=1cm, anchor=west] at (-0.8,-1.6) {cr lf};
  \node[f, fill=lphy!60, minimum width=7.4cm, minimum height=1cm, anchor=north west] at (-0.8,-2.0) {corpo: l'oggetto richiesto};
  \draw[decorate, decoration={brace, amplitude=5pt}] (-1.1,-0.3) -- (-1.1,0.3) node[midway, left=6pt] {riga di stato};
\end{tikzpicture}
\caption{Formato generale di un messaggio di risposta HTTP.}
\end{figure}

\begin{itemize}
    \item la \textbf{versione} di HTTP usata dal server;
    \item il \textbf{codice di stato}: un numero che indica l'esito;
    \item la \textbf{frase}: una descrizione testuale dell'esito.
\end{itemize}

\subsection{I codici di stato}

\begin{table}[H]
\centering
\begin{tabular}{llp{10cm}}
\toprule
\textbf{Classe} & \textbf{Significato} & \textbf{Esempi} \\
\midrule
\texttt{1xx} & informativo & informazioni sulla richiesta in corso \\
\texttt{2xx} & successo & \texttt{200 OK}: richiesta riuscita, l'informazione è nella risposta \\
\texttt{3xx} & redirezione & \texttt{301 Moved Permanently}: l'oggetto è stato spostato in modo permanente; il nuovo URL è
nell'header \texttt{Location} \\
\texttt{4xx} & errore del client & \texttt{400 Bad Request}: il server non ha capito la richiesta;
\texttt{404 Not Found}: il documento non esiste su questo server \\
\texttt{5xx} & errore del server & \texttt{505 HTTP Version Not Supported}: il server non supporta la versione di HTTP richiesta \\
\bottomrule
\end{tabular}
\caption{Le classi dei codici di stato HTTP.}
\end{table}

\info{Come ricordarle}{
    La prima cifra dice di chi è la ``colpa'': \textbf{2} tutto bene, \textbf{3} vai altrove, \textbf{4} hai
    sbagliato tu (client), \textbf{5} ho sbagliato io (server).
}

\subsection{Esempio di risposta}

\begin{lstlisting}[language=http, title={Risposta HTTP}]
HTTP/1.1 200 OK
Connection: close
Date: Tue, 18 Aug 2015 15:44:04 GMT
Server: Apache/2.2.3 (CentOS)
Last-Modified: Tue, 18 Aug 2015 15:11:03 GMT
Content-Length: 6821
Content-Type: text/html

(data data data data data ...)
\end{lstlisting}

La riga di stato dice che il server usa HTTP/1.1 e che è tutto a posto: ha trovato l'oggetto e lo sta
inviando. Il corpo contiene l'oggetto (\texttt{data data \dots}). Gli header:

\begin{itemize}
    \item \texttt{Connection: close}: la connessione verrà chiusa dopo questo messaggio (non persistente);
    \item \texttt{Date}: data e ora in cui la risposta è stata creata e inviata;
    \item \texttt{Server}: il messaggio è stato generato da un server Apache (è l'analogo di \texttt{User-agent});
    \item \texttt{Last-Modified}: data e ora in cui l'oggetto è stato creato o modificato l'ultima volta;
          fondamentale per il caching, perché le copie in cache possono essere vecchie;
    \item \texttt{Content-Length}: il numero di byte dell'oggetto;
    \item \texttt{Content-Type}: il tipo dell'oggetto, qui testo HTML. Il tipo è indicato ufficialmente da
          questo header, \textbf{non} dall'estensione del file.
\end{itemize}

% ══════════════════════════════════════════════════════════════
\section{I cookie}
% ══════════════════════════════════════════════════════════════

Un server HTTP senza stato è semplice, consuma poche risorse e regge migliaia di connessioni TCP
simultanee. Ma la mancanza di stato è anche un grosso limite: molte funzioni del Web dipendono
dal client (il carrello di Amazon, i suggerimenti di Netflix, \dots).

\dfn{Cookie}{
    Un \textbf{cookie} è un gettone digitale (un identificativo alfanumerico) creato dal server e
    consegnato al client, con cui il server \textbf{riconosce un client specifico} nelle richieste
    successive. I cookie possono avere attributi, come una data di scadenza.
}

\subsection{Le quattro componenti}

\begin{enumerate}
    \item Una riga di header \texttt{Set-cookie} nel messaggio di \textbf{risposta} HTTP.
    \item Una riga di header \texttt{Cookie} nel messaggio di \textbf{richiesta} HTTP.
    \item Un \textbf{file} sul sistema del client, gestito dal browser.
    \item Un \textbf{database} lato server (\emph{back-end}).
\end{enumerate}

\begin{figure}[H]
\centering
\begin{tikzpicture}[yscale=0.78, every node/.style={font=\footnotesize\sffamily}]
  \node (cl) at (0,0.9) {\icon[0.9cm]{pc}};
  \node (sv) at (8,0.9) {\icon[0.6cm]{server}};
  \node[below=6pt of cl] {client}; \node[below=6pt of sv] {server Amazon};
  \draw[thick] (0,0) -- (0,-11); \draw[thick] (8,0) -- (8,-11);
  \node[draw, fill=llink!60, rounded corners, anchor=east, text width=2.4cm, align=left, font=\scriptsize\sffamily] at (-0.4,-1) {file dei cookie:\\ \texttt{ebay: 8734}};
  \draw[msg, netblue] (0,-1.2) -- node[above, sloped] {richiesta HTTP normale} (8,-2.2);
  \node[draw, fill=lphy!50, rounded corners, anchor=west, text width=3.1cm, align=left, font=\scriptsize\sffamily] (db) at (8.4,-2.7) {crea l'ID \texttt{1678}\\ e una voce nel database};
  \node[cylinder, shape border rotate=90, aspect=0.25, draw, fill=lphy, minimum height=0.8cm, minimum width=0.7cm] at (12.4,-2.7) {DB};
  \draw[msg, netgreen] (8,-3.2) -- node[above, sloped] {risposta + \texttt{Set-cookie: 1678}} (0,-4.2);
  \node[draw, fill=llink!60, rounded corners, anchor=east, text width=2.4cm, align=left, font=\scriptsize\sffamily] at (-0.4,-4.4) {\texttt{ebay: 8734}\\ \texttt{amazon: 1678}};
  \draw[msg, netblue] (0,-5.4) -- node[above, sloped] {richiesta + \texttt{Cookie: 1678}} (8,-6.4);
  \node[anchor=west, text width=3.6cm, align=left, font=\scriptsize\sffamily] at (8.4,-6.6) {azione specifica per l'utente 1678 (carrello, suggerimenti)};
  \draw[msg, netgreen] (8,-7.0) -- node[above, sloped] {risposta personalizzata} (0,-8.0);
  \node[anchor=east, text=netred, font=\scriptsize\sffamily] at (-0.4,-8.8) {una settimana dopo\dots};
  \draw[msg, netblue] (0,-9.0) -- node[above, sloped] {richiesta + \texttt{Cookie: 1678}} (8,-10.0);
  \node[anchor=west, text width=3.6cm, align=left, font=\scriptsize\sffamily] at (8.4,-10.2) {il server riconosce ancora l'utente};
\end{tikzpicture}
\caption{Un client che usa già eBay contatta Amazon per la prima volta: nasce un cookie, che accompagnerà tutte le richieste successive.}
\end{figure}

\begin{esempio}{Il cookie di Amazon passo per passo}
\begin{enumerate}
    \item Un host, che usa abitualmente eBay (e ha già un cookie per eBay), contatta Amazon per la
          \textbf{prima volta}: non ha alcun cookie per Amazon.
    \item Il server di Amazon crea un \textbf{identificativo} (\texttt{1678}) e una voce corrispondente nel
          proprio database, e risponde includendo l'header \texttt{Set-cookie: 1678}.
    \item Il browser aggiunge al file dei cookie una riga con il nome del server e l'identificativo.
    \item Da quel momento ogni richiesta verso Amazon contiene \texttt{Cookie: 1678}: il server può
          tracciare l'attività dell'utente nel database, sapendo esattamente quali pagine ha visitato, in che
          ordine e quando.
    \item Se l'utente torna una settimana dopo, il browser continua a inserire \texttt{Cookie: 1678}.
\end{enumerate}
\end{esempio}

\subsection{Usi e problemi dei cookie}

I siti usano i cookie per:

\begin{itemize}
    \item il \textbf{carrello}: la lista degli acquisti in corso durante la navigazione;
    \item la \textbf{registrazione}: associando al cookie nome, e-mail, indirizzo, carta di credito, così da
          non doverli reinserire ogni volta;
    \item i \textbf{suggerimenti}: prodotti consigliati in base alle pagine visitate.
\end{itemize}

\warning{Cookie e privacy}{
    I cookie semplificano la vita, ma sono \textbf{controversi}: possono essere considerati un'invasione della
    privacy. Combinando cookie e dati forniti dall'utente, un sito può sapere moltissimo di una persona e,
    potenzialmente, vendere queste informazioni a terzi.
}

% ══════════════════════════════════════════════════════════════
\section{Il caching Web}
% ══════════════════════════════════════════════════════════════

\dfn{Web cache e proxy}{
    Una \textbf{Web cache} è un'entità di rete che soddisfa richieste HTTP \textbf{al posto} del server
    di origine. Ha un proprio disco, su cui conserva copie degli oggetti richiesti di recente.
    Le entità intermedie che lavorano per conto di altri si chiamano \textbf{proxy server}.
}

Un browser può essere configurato in modo che tutte le sue richieste passino prima dalla cache, per
vedere se ha già una copia dell'oggetto.

\begin{figure}[H]
\centering
\begin{tikzpicture}[every node/.style={font=\footnotesize\sffamily}]
  \node (c1) at (0,1.2) {\icon[0.9cm]{pc}}; \node (c2) at (0,-1.2) {\icon[1cm]{laptop}};
  \node[below=-2pt of c2] {client};
  \node (px) at (4.5,0) {\icon[0.7cm]{server}}; \node[below=-1pt of px, align=center] {Web cache\\(proxy)};
  \node (o1) at (9,1.2) {\icon[0.6cm]{server}}; \node (o2) at (9,-1.2) {\icon[0.6cm]{server}};
  \node[below=-1pt of o2] {server di origine};
  \draw[msg, netblue] (c1) -- node[above, sloped] {1. richiesta} (px);
  \draw[msg, netgreen, dashed] (px) to[bend left=15] node[below, sloped] {4. risposta} (c1);
  \draw[msg, netblue] (px) -- node[above, sloped] {2. richiesta (se manca)} (o1);
  \draw[msg, netgreen, dashed] (o1) to[bend left=15] node[below, sloped] {3. risposta} (px);
  \draw[msg, netblue, {Stealth}-{Stealth}] (c2) -- (px);
  \draw[msg, netblue, {Stealth}-{Stealth}] (px) -- (o2);
\end{tikzpicture}
\caption{La cache sta tra i client e i server di origine: è server verso i client e client verso i server.}
\end{figure}

\subsection{Il funzionamento}

Supponiamo che un browser chieda \url{http://www.someschool.edu/campus.gif} passando per una cache:

\begin{enumerate}
    \item Il browser apre una connessione TCP con la \textbf{cache} e le invia la richiesta HTTP.
    \item La cache controlla se ha una copia locale dell'oggetto. Se sì, la restituisce subito al browser
          in una risposta HTTP.
    \item Se no, la cache apre una connessione TCP con il \textbf{server di origine}
          (\texttt{www.someschool.edu}) e gli chiede l'oggetto.
    \item Ricevuto l'oggetto, la cache ne \textbf{salva una copia} e ne invia un'altra al browser.
\end{enumerate}

\warning{La cache è client e server insieme}{
    Una cache è contemporaneamente un \textbf{server} (quando fornisce oggetti ai browser) e un
    \textbf{client} (quando li chiede ai server di origine).
}

\subsection{Perché usare le cache}

\begin{enumerate}
    \item Una cache può \textbf{ridurre molto il tempo di risposta}, soprattutto se tra il client e la cache
          c'è un collegamento veloce e tra la cache e Internet uno lento.
    \item Una cache può \textbf{ridurre molto il traffico} di un'azienda o di un ente verso Internet, e quindi
          il costo della banda.
\end{enumerate}

\begin{figure}[H]
\centering
\begin{minipage}{0.34\linewidth}\centering
  \includegraphics[width=0.9\linewidth]{cap05/cache_istituzionale.png}
\end{minipage}\hfill
\begin{minipage}{0.62\linewidth}
Le cache vengono spesso installate nella rete locale di un'azienda o di un ente.
\dfn{Hit, miss e hit rate}{
    Si ha un \textbf{hit} quando la cache fornisce un oggetto \textbf{senza contattare} il server di
    origine; un \textbf{miss} quando deve andarlo a prendere. La \textbf{hit rate} è la frazione di
    richieste soddisfatte dalla cache: tipicamente \textbf{tra 0{,}2 e 0{,}7}, e cresce con il numero di
    client che usano la cache.
}
Fino al 70\% delle richieste, quindi, può essere servito localmente.
\end{minipage}
\caption{Una cache istituzionale: il percorso rosso (miss) esce su Internet, quello verde (hit) resta nella rete locale.}
\end{figure}

\subsection{Il GET condizionale}

Il caching introduce un problema nuovo: la copia in cache potrebbe essere \textbf{vecchia}, se nel
frattempo l'oggetto è cambiato sul server.

\dfn{GET condizionale}{
    Un \textbf{GET condizionale} è una richiesta HTTP con metodo GET e header
    \texttt{If-Modified-Since}. Il server invia l'oggetto \textbf{solo se} è stato modificato dopo la data
    indicata; altrimenti risponde \texttt{304 Not Modified}, senza corpo.
}

\begin{figure}[H]
\centering
\begin{tikzpicture}[yscale=0.8, every node/.style={font=\scriptsize\sffamily}]
  \timeline{a}{0}{Host A}{9}
  \timeline{c}{5}{Cache}{9}
  \timeline{s}{10}{Server}{9}
  \draw[msg, netblue] (0,-0.4) -- node[above, sloped] {GET oggetto} (5,-1.0);
  \draw[msg, netblue] (5,-1.1) -- node[above, sloped] {GET oggetto} (10,-1.7);
  \draw[msg, netgreen] (10,-1.9) -- node[above, sloped] {200 OK, Last-Modified: X, oggetto} (5,-2.5);
  \node[draw, fill=llink!60, rounded corners, font=\scriptsize\sffamily] at (5,-2.95) {salva l'oggetto con data X};
  \draw[msg, netgreen] (5,-3.4) -- node[above, sloped] {200 OK, oggetto} (0,-4.0);
  \node[anchor=east, text=netred] at (-0.2,-4.9) {più tardi, Host B:};
  \draw[msg, netblue] (0,-5.0) -- node[above, sloped] {GET oggetto} (5,-5.6);
  \draw[msg, netred] (5,-5.8) -- node[above, sloped] {GET oggetto, If-Modified-Since: X} (10,-6.4);
  \draw[msg, netred] (10,-6.6) -- node[above, sloped] {304 Not Modified (nessun corpo)} (5,-7.2);
  \draw[msg, netgreen] (5,-7.5) -- node[above, sloped] {200 OK, oggetto (dalla cache)} (0,-8.1);
\end{tikzpicture}
\caption{La cache verifica con un GET condizionale che la propria copia sia ancora valida: se lo è, l'oggetto non viaggia di nuovo.}
\end{figure}

Se l'oggetto è aggiornato, alla cache basta una piccola risposta del server (niente oggetto da
ritrasmettere) e il client riceve la copia locale.

\subsection{Il caching nel browser}

Anche i \textbf{browser} fanno caching in locale, con lo stesso principio: conservano gli oggetti già
scaricati per non doverli richiedere al server. È una tecnica comunissima che migliora drasticamente
le prestazioni. La copia locale però può non essere aggiornata, e questo genera errori piuttosto
frequenti (la classica pagina che ``non si aggiorna'' finché non si svuota la cache).

\subsection{Oltre il caching: i proxy istituzionali}

\begin{figure}[H]
\centering
\begin{tikzpicture}[every node/.style={font=\footnotesize\sffamily}]
  \node (c) at (0,0) {\icon[1cm]{pc}}; \node[below=-2pt of c] {client};
  \node[netcloud, minimum width=4.2cm, minimum height=2.2cm] (inst) at (4.6,0) {};
  \node at (4.0,0.35) {rete istituzionale};
  \node (px) at (5.3,-0.3) {\icon[0.55cm]{server}}; \node[font=\scriptsize\sffamily] at (5.3,-0.95) {proxy};
  \node[netcloud, minimum width=2.6cm] (inet) at (9.2,0) {Internet};
  \node (o) at (12.1,0) {\icon[0.7cm]{server}}; \node[below=-1pt of o] {server di origine};
  \draw[msg, netblue] (c) -- (px);
  \draw[msg, netblue] (px) -- (inet); \draw[msg, netblue] (inet) -- (o);
  \node[text width=6cm, align=center, text=netred] at (8.2,-1.8) {per il server di origine la richiesta arriva dall'istituzione, non dal client};
\end{tikzpicture}
\caption{Un proxy istituzionale: le richieste escono a nome dell'ente.}
\end{figure}

Oltre alle prestazioni, i proxy servono ad \textbf{accedere a servizi istituzionali}. La rete
dell'Ateneo, per esempio, offre un proxy Web (\texttt{proxy.unina.it}): poiché le richieste vengono
inoltrate dal proxy, per il server di origine \textbf{arrivano tutte dall'istituzione}. Così, da casa, si
possono consultare riviste scientifiche abbonate dall'Università. Esistono anche molti proxy
commerciali o gratuiti su Internet.

% ══════════════════════════════════════════════════════════════
\section{Riepilogo}
% ══════════════════════════════════════════════════════════════

\riepilogo{
\begin{itemize}
    \item Richieste e risposte HTTP sono testo ASCII: \textbf{riga di richiesta/stato}, \textbf{header},
          \textbf{riga vuota}, \textbf{corpo}.
    \item Metodi principali: \textbf{GET}, \textbf{POST}, \textbf{HEAD}, \textbf{PUT}, \textbf{DELETE}.
    \item Codici di stato: 1xx informativi, 2xx successo, 3xx redirezione, 4xx errore del client, 5xx errore del server.
    \item I \textbf{cookie} (Set-cookie / Cookie + file sul client + database sul server) danno uno stato a un
          protocollo stateless, con implicazioni sulla privacy.
    \item Le \textbf{Web cache} (proxy) riducono tempi e traffico; il \textbf{GET condizionale}
          (\texttt{If-Modified-Since} / \texttt{304 Not Modified}) evita copie vecchie.
\end{itemize}
}
